\documentclass[11pt,a4paper]{article}
\usepackage[utf8]{inputenc}
\usepackage{geometry}
\geometry{margin=1in}
\usepackage{graphicx}
\usepackage{hyperref}
\usepackage{enumitem}
\usepackage{indentfirst}
\usepackage{mathpazo}

\begin{document}

\centerline{\Huge{Worcester Polytechnic Institute}}
\vspace{2mm}
\centerline{\huge{Student Branch of the}}
\vspace{4mm}
\centerline{\Huge{Institute of 
Electrical and Electronics Engineers}}
\vspace{2mm}
\centerline{{\fontsize{30}{36}\selectfont Branch Bylaws}}
\hline
\vspace{3mm}
\centerline{\huge{Approved September 30th 2026}}
\bigskip \bigskip \bigskip \bigskip
\bigskip \bigskip \bigskip \bigskip
\bigskip \bigskip \bigskip \bigskip
\begin{figure}[h]
\centering
\includegraphics[width=0.8\linewidth]{Screenshot 2026-09-10 140759.png}
\end{figure}
\newpage
\section*{Preamble}

The Worcester Polytechnic Institute (WPI) Institute of Electrical and Electronics Engineers (IEEE) Student Branch, like its parent organization, was founded with the purpose of "advancing technology for the betterment of humanity". To this end, the WPI IEEE Student Branch serves to enrich the collegiate experience of WPI's students, faculty, and staff who have an interest in electronics and technology.

This document shall outline the Bylaws of the WPI IEEE Student Branch.

\tableofcontents
\vspace{2em}

\newpage
\section*{Bylaw A - Basic Laws}
\addcontentsline{toc}{section}{Bylaw A: Basic Laws}

\subsection*{A.1 Name}
The name of this organization shall be the Institute of Electrical and Electronics Engineers (IEEE) Worcester Polytechnic Institute (WPI) Student Branch, hereafter referred to as the Branch.

\subsection*{A.2 External Policy}
The Branch shall abide by all policies put forth by IEEE and WPI, including, but not limited to, anti-discrimination policies and anti-hazing policies. Furthermore, the Branch shall abide by all federal, state, and local laws.

\subsection*{A.3 Parliamentary Procedure}
All Branch proceedings will follow Robert's Rules of Order for guidance on parliamentary rulings and procedures.

\subsection*{A.4 Faculty Advisor}
The Branch shall be advised by a WPI faculty member, hereafter referred to as the faculty advisor. The responsibilities and appointment of the faculty advisor shall be based upon IEEE's recommendation, which follows. The faculty advisor shall be appointed by the IEEE Regional Director upon the recommendation of the executive board and the WPI Electrical and Computer Engineering Department.

\section*{Bylaw B - Updating the Bylaws}
\addcontentsline{toc}{section}{Bylaw B: Updating the Bylaws}

\subsection*{B.1 Schedule}
To keep the Bylaws up to date and in line with the Branch's needs, the Bylaws are to be reviewed and updated, if found necessary, once every two years following the process described in B.2

\subsection*{B.2 Bylaws Update Procedure}
To update the Bylaws, the Chair and Vice Chair must meet to discuss and draft possible changes. When all the desired updates have been drafted, the Chair and Vice Chair will present these changes to the executive board at an executive board meeting. Each change, individually, will then be discussed and voted on as a formal motion (See \hyperref[F.2]{F.2}).

\section*{Bylaw C - Membership}
\addcontentsline{toc}{section}{Bylaw C: Membership}

\subsection*{C.1 Classes of Membership}
The Branch shall have three classes of membership:
\begin{itemize}[noitemsep]
    \item \textbf{General Member:} A WPI student who has attended at least one Branch event, but does not satisfy the requirements of being an active member.
    \item \textbf{Active Member:} A WPI student who satisfies the requirements of active member status, as defined in \hyperref[C.2]{C.2}.
    \item \textbf{Executive:} A WPI undergraduate student who is a member of the executive board.
\end{itemize}

\subsection*{C.2 Requirements of 
Active Member Status}
\label{C.2}
Active member status shall be granted to WPI students who have at least 3 membership points. This threshold assumes that Branch events are held approximately every week throughout the term, for all four terms in the academic year. Adjust this threshold if this event frequency is not met. Point assignment values are specified in \hyperref[E.2]{E.2}.

Membership points are accumulated across a rolling semester. A rolling semester is defined as a contiguous span of two terms where a student is not on leave. If a student is on leave, they must notify the executive board to exclude the term they are away from the rolling period. The executive board will evaluate exceptions under extenuating circumstances on a case-by-case basis.

\section*{Bylaw D - Executive Board}
\addcontentsline{toc}{section}{Bylaw D: Executive Board}

\subsection*{D.1 Components}
The executive board shall serve as the central authority of the Branch. The executive board shall be composed of the following positions, hereafter referred to as executives:
\begin{enumerate}[noitemsep]
    \item Chair
    \item Vice Chair
    \item Treasurer
    \item Secretary
    \item Events Chairs
    \item Projects Chair
    \item Public Relations Chair
    \item Web Chair
\end{enumerate}
Executives shall be elected as outlined in \hyperref[G]{Bylaw G}.
When "highest-ranking executive" is referenced, use the order of the list above as the ranking.

\subsection*{D.2 Expectations of Executives}
Executives shall be expected to:
\begin{itemize}[noitemsep]
    \item Maintain the status of an active member.
    \item Be in attendance, excused, or on leave for all Branch activities, including executive board meetings.
    \item Satisfy all requirements specific to the executive's position.
    \item Share the responsibilities of the executive board.
    \item Maintain familiarity with the Branch bylaws.
    \item Maintain IEEE Student Membership throughout the duration of their term as an executive.
\end{itemize}


\subsection*{D.3 25 Live Access Priority}
25 Live is WPI’s platform for booking campus spaces. Because the number of authorized users fluctuates and is determined by the Student Activities Office (SAO), the following list establishes the priority order for receiving access:
\begin{enumerate}[noitemsep]
    \item Chair
    \item Events Chairs
    \item Projects Chair
    \item Vice Chair
    \item Secretary
    \item Treasurer
    \item Public Relations Chair
    \item Web Chair
\end{enumerate}


\subsection*{D.4 Chair}
The Chair may also be referred to as the President in accordance with WPI Executive Board regulations. It shall be the responsibility of the Chair to oversee the operations of the Branch in its entirety. This duty shall entail:
\begin{itemize}[noitemsep]
    \item Meeting with executives to discuss and develop goals and achievements.
    \item Moderating executive board meetings.
    \item Assigning tasks to executives that fall outside their designated responsibilities.
    \item Preparing an annual Branch plan during A-Term and submitting it in vTools.
    \item Submitting at least four event reports in vTools over the academic year.
    \item Maintaining executive reporting in vTools.
    \item Serving as a representative on the Department of Electrical and Computer Engineering (ECE) Student Advisory Board under the ECE Department Head.
    \item Maintaining the Branch Bylaws.
\end{itemize}

\subsection*{D.5 Vice Chair}
The Vice Chair may also be referred to as the Vice President in accordance with WPI Executive Board regulations. It shall be the responsibility of the Vice Chair to oversee the internal operations of the Branch. This duty shall entail:
\begin{itemize}[noitemsep]
    \item Assisting the Chair with their duties.
    \item Helping fellow executives discuss and develop goals and achievements.
    \item Holding fellow executives accountable for failing to meet responsibilities.
    \item Providing fellow executives with resources and assistance as needed.
    \item Upholding the Bylaws.
    \item Ensuring strong coordination and cooperation amongst all executives.
\end{itemize}

\subsection*{D.6 Treasurer}
It shall be the responsibility of the treasurer to oversee the Branch finances. Although they are the sole member with access to funds, they are not the sole moderator of the funds' usage, but rather a guide for how the funds are to be used. This duty shall entail:
\begin{itemize}[noitemsep]
    \item Serving as the financial advisor for the Branch.
    \item Submitting an annual budget proposal to the WPI Student Government Association (SGA).
    \item Submitting Funding Requests to the SGA a minimum of 2 weeks before the intended use date.
    \item Pursuing funds from external sources including, but not limited to, the IEEE Worcester County Section (WCS), the WPI ECE department, and corporate sponsors.
    \item Maintaining Branch funds:
        \begin{itemize}[noitemsep]
            \item Within the Branch's budget.
            \item Received from SGA via funding requests.
            \item Received from outside sponsors.
            \item Within the Branch's gift fund held by the ECE department.
        \end{itemize}
    \item Maintaining a history of Branch finances.
    \item Completing reimbursements for pre-approved purchases.
    
\end{itemize}

\subsection*{D.7 Secretary}
It shall be the responsibility of the secretary to oversee the administration of the Branch. This duty shall entail:
\begin{itemize}[noitemsep]
    \item Recording meeting minutes during executive board meetings and sending out action items afterward.
    \item Maintaining member attendance records at Branch activities.
    \item Updating member activity status as required.
    \item Maintaining the Branch file system, which includes the Google Drive folder.
    \item Creating and sending out the Branch weekly newsletter.
    \item Ensuring synchronization between myWPI members and the email alias.
\end{itemize}

\subsection*{D.8 Events Chairs}
There shall be two Events Chairs, and it will be the responsibility of the Events Chairs to coordinate all logistics and planning for Branch activities, including coordinating with other executives when needed. This duty shall entail:
\begin{itemize}[noitemsep]
    \item Arranging weekly events, which must include the following each term of the academic year:
    \begin{itemize}[noitemsep]
        \item One General Body Meeting (GBM)
        \item One technical event
        \item One social event
        \item One flagship event
    \end{itemize}
    \item Uploading event details into MyWPI 1.5-2 weeks before said event takes place.
    \item Filling out the event template 1.5-2 weeks before said event takes place.
    \item Being the main point of contact for other student groups on campus for collaboration.
    \item Maintaining the events calendar.
\end{itemize}

\subsection*{D.9 Projects Chair}
It will be the responsibility of the Projects Chair to create and maintain technical projects that promote and/or benefit the Branch and WPI community. This duty shall entail:
\begin{itemize}[noitemsep]
    \item Identifying at least one project that can be completed or improved upon every academic year.
    \item Maintaining existing projects to ensure that they are operational.
    \item Open-sourcing and maintaining all projects on the Branch's GitHub organization.
    \item Managing the Projects budget.
\end{itemize}

\subsection*{D.10 Public Relations Chair}
It will be the responsibility of the Public Relations Chair to promote and document Branch events. This duty shall entail:
\begin{itemize}[noitemsep]
    \item Designing graphics for events.
    \item Promoting events and meetings in ways including, but not limited to, posting to social media, hanging posters and fliers, and advertising on campus televisions.
    \item Documenting events by taking photos and/or videos and posting recaps.
    \item Coming up with ideas for WPI IEEE Branded Merchandise.
\end{itemize}

\subsection*{D.11 Web Chair}
It will be the responsibility of the Web Chair to maintain, update, and improve the Branch website. This duty shall entail:
\begin{itemize}[noitemsep]
    \item Ensuring the website is publicly accessible at all times.
    \item Ensuring the website's code base is organized and version controlled.
    \item Making sure events, executives, projects, and other content on the website is up to date.
    \item Maintaining the website calendar.
    \item Improving the website by adding or improving features each academic term.
\end{itemize}


\section*{Bylaw E - Meetings}
\addcontentsline{toc}{section}{Bylaw E: Meetings}

\subsection*{E.1 Classes of Meetings}
\label{E.1}
The Branch shall have four classes of meetings:
\begin{itemize}[noitemsep]
    \item \textbf{General Body Meeting:} A meeting held at least once per term, ideally at the beginning of the term, that is open to all WPI students. During general body meetings, the branch will discuss updates and future events. Elections may also be run during or in place of the C-Term General Body Meeting.
    \item \textbf{Events:} Organized activities, gatherings, or initiatives held to foster engagement, promote the club's mission, and provide opportunities for member involvement. Events may include workshops, social gatherings, professional development opportunities, and other activities aligned with the club's objectives.
    \item \textbf{Flagship Event:}
    Flagship events are large-scale events that happen once per academic term. The current set of flagship events is described below. Each flagship event requires the formation of a committee to assist in the planning and execution of these large events (See \hyperref[I]{Bylaw I} for a description of committees).
    \begin{itemize}[noitemsep]
        \item \textbf{A-Term: PCB Design Class} -- A project-oriented hardware design course where students learn PCB design over several weeks, then design, receive, and assemble their own PCBs. (Location: Flexible, Date: Early September)
        \item \textbf{B-Term: Spark Party} -- A night to celebrate the ECE department's accomplishments, featuring music and dance performances from student groups, WCS demonstrations, faculty and MQP presentations, with a finale tesla coil performance. (Location: AK 116, Date: Early December)
        \item \textbf{C-Term: Networking Night} -- A unique opportunity for highly motivated, talented, technical WPI students across different majors to connect with top companies in the industry over a catered dinner. (Location: CC Odeum, Date: Late January)
       % \item \textbf{D-Term: Hackathon} -- A weekend-long project-focused hackathon with ample hardware and facilities open to students over the weekend. (Location: Unity Hall, Date: Late March)
    \end{itemize}
    \item \textbf{Executive Board Meeting:} A 1-hour meeting held once per week during all academic terms that is open to the WPI community. During executive board meetings, the executive board discusses and plans Branch activities and votes on Branch matters.
\end{itemize}

\subsection*{E.2 Meeting Attendance}
\label{E.2}
There shall be four classes of meeting attendance for all Branch meetings:
\begin{itemize}[noitemsep]
    \item \textbf{In Attendance:} The person was present for the majority of the meeting.
    \item \textbf{Excused:} The person was not present for the majority of the meeting, and they sent a valid excuse to the executive board prior to the start of the meeting.
    \item \textbf{Absent:} The person was not present for the majority of the meeting, and they did not send a valid excuse to the executive board prior to the start of the meeting.
    \item \textbf{On Leave:} The person was away from campus for the academic term in which the meeting was held for reasons including but not limited to Interactive Qualifying Project (IQP), Major Qualifying Project (MQP), or Co-op.
\end{itemize}
Branch activities will be assigned the following point values:
\begin{itemize}[noitemsep]
    \item Events: 1.0 point
    \item General Body Meetings: 1.0 point
    \item Executive Board Meetings: 0.50 points
    \item PCB Design Lecture: 0.5
\end{itemize}

\section*{Bylaw F - Voting Procedure}
\addcontentsline{toc}{section}{Bylaw F: Voting Procedure}

\subsection*{F.1 Voting Membership}
The voting membership is the set of people who are eligible to vote in a certain ballot. For executive elections, the voting membership is the set of active members. For formal motions at executive board meetings, the voting membership is the set of executives.

\subsection*{F.2 Formal Motions}
\label{F.2}
During an executive board meeting, any executive may make a formal motion. The motion shall be discussed and, if necessary, voted on as outlined in Bylaw \hyperref[F.3]{F.3}. Motions include, but are not limited to, the following:
\begin{itemize}[noitemsep]
    \item Amending the Branch Bylaws.
    \item Approving purchases with Branch funds.
    \item Impeaching an executive.
\end{itemize}

\subsection*{F.3 Passing a Motion}
\label{F.3}
When balloting for a motion, each member of the voting membership is entitled to a single vote either in favor of or against the motion. The motion will pass provided that two-thirds of the voting membership votes in favor of the motion. At the discretion of the executive board, the votes may be cast either publicly or anonymously. If the votes are cast publicly, the moderator shall not be entitled to a vote.

\subsection*{F.4 Electing an Executive}
\label{F.4}
Balloting to elect an executive shall occur during a general body meeting and shall proceed as follows:
\begin{enumerate}
    \item If there are no candidates for the position, the position shall be left vacant until an appointment can be made or the duties of the unfilled position have been equitably distributed among the other executives.
    
    \item If there are two candidates for the position, each member of the voting membership shall be entitled to an anonymous vote for one candidate. The candidate with the largest number of votes shall be elected to the executive position. In the case of a tie, current executive members shall hold a ballot to determine the winner. Executive members cannot participate in this ballot if they are one of the tied candidates. If no consensus is reached during this ballot, the current Chair will determine the new executive member.
    
    \item In the event that more than two candidates are nominated for a single position, the election shall be conducted using an anonymous, ranked-choice ballot. Every member of the voting membership shall be given the right to cast a single ballot, on which they must rank up to three candidates in order of their preference (i.e., first choice, second choice, and third choice).
\end{enumerate}

\section*{Bylaw G - Elections}
\addcontentsline{toc}{section}{Bylaw G: Elections}
\label{G}

\subsection*{G.1 Schedule}
The election of the executive board, hereafter referred to as the election, shall be held at least once a year and shall take place no later than the second-to-last full week of C-term.

\subsection*{G.2 Nominations}
The executive board shall notify all active members of the election two weeks prior to the election. Any active member may nominate any undergraduate active member. Active members may be nominated for more than one position. Nominations may be submitted to the executive board up until the start of the election. Upon receiving a nomination, the executive board will immediately notify the nominee. To be a candidate in the election, the nominee must accept the nomination prior to the start of the election. The nomination is assumed to be automatically accepted in the case of a self-nomination.

%% workshop later
\subsection*{G.3 Moderators}
% Elections for each board position shall be moderated by the two highest-ranking executives who do not intend to run for any elected position. In the event that no such members are available, the current Executive Board shall oversee the election of the Chair. Once a Chair has been elected by the general body, the newly elected Chair shall assume responsibility for moderating the elections for the remaining board positions.

Elections for each board position shall be moderated by the two highest-ranking executives who do not intend to run for any elected position. To maintain impartiality, moderators may not participate in any discussions or deliberations regarding the election of an executive. In the event that no such members are available, the current Executive Board shall oversee the election of the Chair. Once a Chair has been elected by the general body, the newly elected Chair shall assume responsibility for moderating the election for the Vice Chair, who will then join the newly elected Chair in moderating the elections for the remaining board positions.

\subsection*{G.4 Procedure}
When electing an executive the following procedure will be followed:
\begin{enumerate}
    \item The names of all candidates for the position shall be read.
    \item All candidates shall be asked to leave the room.
    \item Candidates shall be selected one by one, in a random order. Each selected candidate shall return to the room and may deliver a speech no longer than three minutes. The general body may then ask the selected candidate questions for up to five minutes, following which the selected candidate shall be asked to leave the room.
    \item Following this, the general body may discuss the candidates for up to 10 minutes.
    \item All candidates shall return to the room, and balloting to elect an executive shall occur as described in Bylaw \hyperref[F.4]{F.4}.
\end{enumerate}

\subsection*{G.5 Transition}
% The newly elected executive board shall take office immediately following D-term. Throughout D-term, the newly elected executive board shall be mentored by the current executive board. This shall include, at minimum, the following:

% Each current executive shall review and/or update their respective positions' transition documents.

% Each newly elected executive shall:
% \begin{enumerate}
%     \item Read the transition documents for their position.
%     \item Meet with the current executive for their position.
%     % reword
%     \item Attend all remaining executive board meetings before taking office in D-term.
%     \item The newly elected executive board shall attend all general body meetings.
% \end{enumerate}

Following elections, the current executive board will mentor the incoming officers throughout C-Term. The newly elected board will officially assume their roles at the start of D-Term. This transition period shall include, at minimum, the following:

Current executives shall review and update the transition documents for their respective positions.

Newly elected executives shall:
\begin{enumerate}
  \item Review the transition materials for their designated roles.

  \item Meet with the outgoing executive holding their position.

  \item Attend all remaining C-Term executive board meetings.

  \item Attend all general body meetings.
\end{enumerate}

\section*{Bylaw H - Vacant Executive Position}
\addcontentsline{toc}{section}{Bylaw H: Vacant Executive Position}

\subsection*{H.1 Vacancy}
Executive positions may become vacant outside of the standard election schedule. This can happen if the executive resigns, is impeached, no executive is elected to the position, the executive graduates early, or the executive is away from campus.

\subsection*{H.2 Delegation of Responsibilities}
While a board position remains vacant, the Executive Board shall determine an appropriate method for distributing the associated responsibilities among current executives until the position is formally filled.

% Not term based cause it can happen whenever
\subsection*{H.3 Appointment of an Interim Executive}
% During any period in which a board position is vacant, the Executive Board shall make reasonable efforts to appoint an interim executive on a term-by-term basis. This process shall include distributing an application to all active members at the beginning of the term and selecting a suitable candidate for appointment.

During any period in which a board position is vacant, the Executive Board shall make reasonable efforts to appoint an interim executive. This process shall include distributing an application to all active members and selecting a suitable candidate for appointment.

\subsection*{H.4 Interim Positions}
In the event that an interim is appointed to fill a vacant board position, the interim shall assume all responsibilities and duties associated with that position, as outlined in Bylaw D. These responsibilities shall be upheld until the original executive returns from leave or until the next general board elections are conducted.

\section*{Bylaw I - Committees}
\label{I}
\addcontentsline{toc}{section}{Bylaw I: Committees}

A committee is a collection of executives and active members who work together to coordinate events or reach club goals.

\subsection*{I.1 Creating a Committee}
Committees are required for each of the flagship events (See \hyperref[E.1]{E.1}), as these events require a lot of work. Other committees may be formed for other tasks (e.g., coordinating going to an IEEE conference) by presenting the need at a board meeting and passing a formal motion (See \hyperref[F.2]{F.2}).

\subsection*{I.2 Committee Duties}
Committees are expected to meet at least once every other week leading up to the committee's respective event or deadline for the committee's goal. Committees must also record the happenings of their meetings in a minutes format to be submitted to the executive board so they stay updated.

\subsection*{I.3 Committee Structure}
All committees must have at least one executive who shall serve as the "Committee Lead". It is the duty of the Committee Lead to make sure the meeting notes are taken and to report back to the executive board. In the event there is more than one executive on a committee, the executives may choose among themselves who is to serve as the Committee Lead, or the title falls to the highest-ranking executive.

% "Ideally composed"
\subsection*{I.4 Spark Party Committee}
% Unlike other committees, the Spark Party Committee is to be composed of 2 members of the executive board, 2 Women in Electrical and Computer Engineering Club (WECE) members, and 2 IEEE Eta Kappa Nu (HKN) members. If a potential committee member is in two or more of the clubs listed, they must choose one club to represent to have no less than 6 members in the committee.

The Spark Party Committee should consist of at least six unique members, ideally comprised of two IEEE Executive Board members, two representatives from the Women in Electrical and Computer Engineering (WECE) club, and two from IEEE Eta Kappa Nu (HKN). Members affiliated with multiple listed organizations must select a single group to represent to prevent overlap and maintain the minimum six-member requirement.

\section*{Bylaw J - Branch Lounge}
\addcontentsline{toc}{section}{Bylaw J: Branch Lounge}

\subsection*{J.1 Lounge Access}
Lounge access is restricted to active members who possess an IEEE membership. This is because IEEE funds the room.

\subsection*{J.2 Lounge Policy}
Below is the Lounge Policy. This is the governing document for conduct in the lounge. This is to be posted in the lounge for reference and to be adhered to. Failure to abide by the Lounge Policy will result in losing lounge access.\newline

By gaining access to the Branch's Lounge in AK 106, you agree to abide by the following rules. Having our own space on campus is a privilege, and these rules help keep the space functional and clean for everyone! The executive board reserves the right to revoke lounge access if these policies are not respected.

\begin{enumerate}
    \item{Regular access to the lounge is reserved exclusively for lounge members. The requirements for lounge members are defined in the bylaws.}
    \item{Guests are permitted for short periods of time under supervision by a lounge member.}
    \item{Usage of lounge space, tools, and resources is first come first serve.}
    \item{Do not host meetings in the lounge unless they are IEEE-related.}
    \item{Do not take other people's belongings. This includes, but is not limited to, food, electronics, and personal items.}
    \item{Do not remove shared items from the lounge. This includes, but is not limited to, supplies, furniture, and equipment.}
    \item{Clean up after you leave. This includes, but is not limited to, ensuring surfaces and project areas are kept tidy.}
    \item{Please try to label food in the fridge with your name. Food over a week old will be thrown out without notice.}
    \item{Be nice to others! We're all friends here!}
\end{enumerate}

\noindent We're trying to build an open, fun, collaborative place on campus. This space belongs to all of us — we want you to bring your personal touch and help make it better! There aren't any rules on what items should/shouldn't be in the lounge, so feel free to bring whatever you want. If there's something you think the club should purchase or restock, please let the executive board know!

\newpage
\section*{Glossary}
\addcontentsline{toc}{section}{Glossary}
\begin{description}[style=nextline, leftmargin=1.5em]
    \item[25Live]
    WPI's campus space reservation platform used by authorized members of the Branch to reserve rooms and other campus spaces for Branch activities.

    \item[Active Member]
    A WPI student who satisfies the active member requirements established in Bylaw C.2, including the required number of membership points within the applicable rolling semester.

    \item[Appointment]
    The selection of an individual by the Executive Board to temporarily fill a vacant executive position outside of the standard election process, as provided in Bylaw H.

    \item[Branch]
    The Institute of Electrical and Electronics Engineers Worcester Polytechnic Institute Student Branch, also referred to as the WPI IEEE Student Branch.
    
    \item[Branch Activity]
    Any meeting, event, project, committee activity, program, or other activity organized or officially recognized by the Branch.

    \item[Branch Bylaws]
    This document and its provisions governing the organization, membership, Executive Board, meetings, elections, committees, and other operations of the Branch.

    \item[Branch Funds]
    Money available to or administered by the Branch, including its annual budget, funding received through SGA funding requests, external sponsorships, and funds maintained through the Branch's ECE Department gift fund.
    
    \item[Candidate]
    An eligible nominee who has accepted a nomination for an elected executive position.

    \item[Committee Lead]
    The executive responsible for overseeing a committee, ensuring that committee meeting notes are recorded, and reporting committee activity to the Executive Board.

    \item[ECE]
    The Worcester Polytechnic Institute Department of Electrical and Computer Engineering.

    \item[Executive]
    A WPI undergraduate student who holds a position on the Executive Board.

    \item[Executive Board]
    The central governing authority of the Branch, composed of the Chair, Vice Chair, Treasurer, Secretary, Events Chairs, Projects Chair, Public Relations Chair, and Web Chair.
    
    \item[Executive Board Meeting]
    A meeting normally held once each week during WPI academic terms at which the Executive Board discusses, plans, and votes on Branch matters. Executive Board Meetings are open to the WPI community.
    
    \item[Executive Term]
    The period during which an elected executive officially holds their Executive Board position.

    \item[Faculty Advisor]
    A WPI faculty member appointed to advise the Branch in accordance with IEEE requirements and Bylaw A.4.

    \item[Flagship Event]
    A major Branch event intended to occur once per academic term and requiring the formation of a committee to assist with its planning and execution.

    \item[Formal Motion]
    A formal proposal made by an executive during an Executive Board Meeting that is presented for discussion and, when necessary, a vote in accordance with Bylaw F.
    
    \item[Funding Request]
    A request submitted to the WPI Student Government Association for funding to support a Branch expense or activity.

    \item[General Body]
    The broader membership participating in Branch activities. During elections or other matters involving voting, only individuals who meet the applicable definition of the voting membership may cast a ballot.

    \item[GitHub Organization]
    The Branch's organizational account on GitHub used to host, version-control, and maintain Branch software and technical projects.

    \item[Google Drive]
    The Branch's electronic file storage system maintained as part of the Branch's administrative records.
    
    \item[Highest-Ranking Executive]
    The eligible executive holding the highest position according to the ranking established in Bylaw D.1. The ranking, from highest to lowest, is Chair, Vice Chair, Treasurer, Secretary, Events Chairs, Projects Chair, Public Relations Chair, and Web Chair.

    \item[HKN]
    IEEE Eta Kappa Nu, the honor society of the ECE department. Also incorporated with IEEE.

    \item[IEEE]
    The Institute of Electrical and Electronics Engineers.

    \item[IEEE Student Membership]
    Student-level membership in IEEE. Executives are required to maintain IEEE Student Membership throughout their term in office.

    \item[Impeachment]
    The process of removing an executive from office through a formal motion and vote in accordance with the applicable voting requirements of these Bylaws.

    \item[Interim Executive]
    An individual appointed by the Executive Board to temporarily assume the responsibilities and duties of a vacant executive position until the original executive returns from leave or the position is filled through the next general election.

    \item[IQP]
    Interactive Qualifying Project, a WPI degree requirement that may involve a student being away from campus and therefore considered On Leave for purposes of these Bylaws.

    \item[Leave / On Leave]
    A meeting and membership status applying to a person who is away from campus for an academic term for reasons including, but not limited to, an IQP, MQP, or Co-op. A student seeking to exclude such a term from their rolling semester must notify the Executive Board.

    \item[Lounge]
    The Branch space located in AK 106 and governed by Bylaw J and the Lounge Policy.

    \item[Lounge Access]
    Authorization to regularly use the Branch Lounge. Under Bylaw J, regular access is restricted to active members who also possess IEEE membership.

    \item[Lounge Member]
    An active member who possesses IEEE membership and is therefore eligible for regular access to the Branch Lounge.

    \item[Majority of a Meeting]
    More than half of the duration of a meeting. Attendance for the majority of a meeting constitutes the attendance status "In Attendance."

    \item[Membership Point]
    A unit awarded for participation in qualifying Branch activities and used to determine active member status in accordance with Bylaws C.2 and E.2.

    \item[Moderator]
    An executive designated to oversee the procedure of a meeting, election, discussion, or vote. Election moderators are subject to the impartiality requirements established in Bylaw G.3.

    \item[Motion]
    See \textbf{Formal Motion}.

    \item[MQP]
    Major Qualifying Project, a WPI degree requirement that may involve a student being away from campus and therefore considered On Leave for purposes of these Bylaws.
    
    \item[myWPI]
    WPI's student organization platform used by the Branch to manage organizational information, membership, and event information.

    \item[Nomination]
    The process by which an active member proposes an eligible undergraduate active member for election to an executive position.

    \item[Nominee]
    An eligible active member who has been nominated for an executive position. A nominee becomes a candidate upon accepting the nomination.

    \item[PCB]
    Printed Circuit Board.

    \item[PCB Design Lecture]
    A lecture associated with the PCB Design Class for which attendance may earn membership points as specified in Bylaw E.2.

    \item[President]
    An alternative title for the Chair used where required or recognized under WPI Executive Board regulations.

    \item[Ranked-Choice Ballot]
    The election ballot used when more than two candidates seek a single executive position. Each voter may rank up to three candidates in order of preference.

    \item[Reimbursement]
    Repayment from Branch funds for an authorized purchase that was approved before the purchase was made.
    
    \item[Rolling Semester]
    A contiguous span of two WPI academic terms during which a student is not On Leave. A term of approved leave may be excluded from the rolling period after the student notifies the Executive Board.

    \item[SAO]
    The WPI Student Activities Office.

    \item[SGA]
    The Worcester Polytechnic Institute Student Government Association.

    \item[Vacancy / Vacant Executive Position]
    An Executive Board position that is not currently occupied by its elected executive. A vacancy may result from resignation, impeachment, failure to elect an executive, early graduation, leave, or another circumstance described in Bylaw H.

    \item[Vice President]
    An alternative title for the Vice Chair used where required or recognized under WPI Executive Board regulations.

    \item[vTools]
    IEEE's online organizational reporting system used by the Branch for tasks including Branch planning, activity reporting, and executive reporting.

    \item[WCS]
    The IEEE Worcester County Section.
    
    \item[WECE]
    Women in Electrical and Computer Engineering, a WPI student organization referenced in connection with the Spark Party Committee.
    
    \item[WPI]
    Worcester Polytechnic Institute.
\end{description}

\end{document}